\chapter{Basi a dimensione infinita}

Abbiamo dimostrato che ogni spazio vettoriale a dimensione finita possiede una base. La dimostrazione da noi effettuata non è, tuttavia, valida
quando si passa a dimensione infinita. \\
Si può recuperare, tramite l'assioma della scelta, il concetto di base che si ha a dimensione finita (ovvero di insieme di vettori linearmente indipendenti massimale) anche
per spazi a dimensione infinita; sebbene sia poco pratica negli spazi infinito-dimensionali che solitamente si incontrano (ovvero spazi di Hilbert separabili). Enunciamo innanzitutto l'assioma della scelta

\begin{axiom}[Assioma della scelta]
    Sia $X$ un insieme non vuoto e $\mathcal{P}^*(X) = \mathcal{P}(X) \setminus \{ \emptyset \}$ l'insieme delle parti non vuote di $X$, allora esiste una funzione, detta funzione della scelta, $f: \mathcal{P}^*(X) \to X$ tale che $f(A) \in A \, \forall A \in \mathcal{P}^*(X)$ 
\end{axiom}

Questo è un assioma particolarmente controverso della matematica moderna, siccome comporta una serie di conseguenze paradossali come il paradosso di Banach-Tarski e la costruzioni di insiemi non misurabili secondo Lebesgue, ovvero l'insieme di Vitali; tuttavia, è necessario per 
garantire la possibilità di costruire un insieme scegliendo un elemento da ciascuno degli insiemi: quando abbiamo una famiglia finita di insiemi, la possibilità di fare ciò è garantita dagli assiomi della teoria ZF\footnote{ZF, che sta per Zermelo-Frankel, è una sigla che si usa per indicare gli assiomi alla base della teoria degli insiemi, che necessiterebbero di un approfondimento
troppo sostanzioso per essere inseriti in una mera appendice. Riassumendo parecchio, questi affermano che esiste l'insieme vuoto, è possibile effettuare l'unione e l'intersezione di insiemi ottenendo nuovamente degli insiemi, esiste l'insieme singoletto, ogni insieme non può contenere sé stesso e così via}; mentre quando abbiamo un numero infinito di insiemi nulla ci assicura, dal punto di vista
assiomatico, la possibilità di farlo. Tale assioma, tuttavia, non crea alcuna contraddizione con gli altri che regolano la teoria insiemistica: è stato, infatti, dimostrato che tale assioma è indipendente da quelli di ZF. \\
E' possibile dimostrare, ma non lo riporteremo, che l'assioma della scelta è equivalente al \emph{lemma di Zorn}

\begin{lemma}[Lemma di Zorn]
    Ogni insieme induttivo ha un elemento massimale
\end{lemma}
Per capire bene che cosa questo lemma rappresenta, dobbiamo chiaramente andare a sviscerare il concetto di insiema induttivo e di elemento massimale, quindi riporto una serie di definizioni necessarie
\begin{definition}[relazione]
    Sia $A, B$ due insiemi, definiamo relazione $\mathcal{R}$ un sottoinsieme di $A \times B$. Due elementi $x \in A$ e $y \in B$ sono messi in relazione da $\mathcal{R}$ se
    $$
        (x, y) \in \mathcal{R}
    $$
    e, in tal caso, si scrive $x \mathcal{R} y$.
\end{definition}
\begin{definition}[relazione d'ordine parziale]
    Sia $A$ un insieme e sia $\leq \subseteq A \times A$ una relazione. Diremo che $\leq$ è una relazione d'ordine se
    \begin{itemize}
        \item $x \leq x \, \forall x \in A$ (proprietà riflessiva);
        \item $x \leq y \wedge y \leq x \implies x = y$ (proprietà antisimmetrica);
        \item $x \leq y \wedge y \leq z \implies x \leq z$ (proprietà transitiva).
    \end{itemize}
\end{definition}
L'esempio più semplice di relazione d'ordine è il $\leq$ che abbiamo su $\mathbb{Q}$, tuttavia non è una relazione d'ordine totale, ovvero soddisfa una quarta proprietà (di totalità) secondo cui $\forall x, y \in A, x \leq y \vee y \leq x$.
\begin{definition}[insieme parzialmente ordinato]
    Un coppia $(A, \leq)$ si dice insieme parzialmente ordinato se $\leq$ è una relazione d'ordine parziale definita su $A$.
\end{definition}
\begin{definition}[catena]
    Sia $(A, \leq)$ un insieme parzialmente ordinato, diremo che $C \in \mathcal{P}(A)$ è una catena se è totalmente ordinato dalla relazione d'ordine.
\end{definition}
\begin{definition}[elemento massimale]
    Sia $(A, \leq)$ un insieme parzialmente ordinato, diremo che $x \in A$ è un elemento massimale se $\not\exists y \in A: y > x$.
\end{definition}
